\section{The direct semiprime family}\label{sec:direct-family}

Throughout Sections \ref{sec:direct-family}--\ref{sec:moving-target}, fix a squarefree integer $b\ge3$ and a finite set $T\subset\N$. All implicit constants may depend on $b$ unless explicitly stated otherwise, and $T$ is fixed before $X\to\infty$.

Set
\[
Z=X^3,
\qquad
\mathcal P=\{p\text{ prime}:X\le p<Z\},
\qquad
\mathcal B=\{q\text{ prime}:Z/2\le q<Z\},
\tag{2.1}
\]
and let $S_b$ denote the set of prime divisors of $b$. Define
\[
\mathcal E_{\mathrm{pair}}
=
\{pq:p,q\in\mathcal P,\ p<q\},
\qquad
\mathcal E_b
=
\{rq:r\in S_b,\ q\in\mathcal B\},
\tag{2.2}
\]
\[
\mathcal E
=
\mathcal E_{\mathrm{pair}}\sqcup\mathcal E_b,
\qquad
L=b\prod_{p\in\mathcal P}p.
\tag{2.3}
\]
For sufficiently large $X$, every prime in $\mathcal P$ exceeds every prime divisor of $b$, and every denominator in \eqref{2.2} exceeds $\max T$.

\begin{lemma}[Arithmetic shape]\label{lem:arithmetic-shape}
For all sufficiently large $X$:
\begin{enumerate}[label=\textup{(\roman*)}]
\item every $e\in\mathcal E$ is the product of two distinct primes;
\item the elements of $\mathcal E$ are pairwise distinct and lie outside $T$;
\item every $e\in\mathcal E$ divides $L$;
\item $L$ is squarefree.
\end{enumerate}
\end{lemma}

\begin{proof}
The two factors of every denominator in \eqref{2.2} are distinct once $X$ exceeds the prime divisors of $b$. Unique factorization separates products of two primes from $\mathcal P$ from products $rq$ with $r\mid b$ and $q\in\mathcal B$, and also proves distinctness inside each family. The choice of $X$ gives avoidance of $T$. Since $b$ is squarefree and disjoint in prime support from $\mathcal P$, the integer $L$ is squarefree and every denominator in \eqref{2.2} divides it.
\end{proof}

\subsection{Reciprocal load and variance}

Write
\[
S_X=\sum_{p\in\mathcal P}\frac1p,
\qquad
U_X=\sum_{p\in\mathcal P}\frac1{p^2}.
\]
The prime number theorem and partial summation give
\[
S_X=\log3+o(1),
\qquad
U_X=O\!\left(\frac1{X\log X}\right).
\tag{2.4}
\]
Consequently
\[
\sum_{e\in\mathcal E_{\mathrm{pair}}}\frac1e
=
\frac12(S_X^2-U_X)
\longrightarrow
\lambda_3:=\frac{(\log3)^2}{2}.
\tag{2.5}
\]
For the target rows,
\[
\sum_{e\in\mathcal E_b}\frac1e
=
\left(\sum_{r\in S_b}\frac1r\right)
\left(\sum_{q\in\mathcal B}\frac1q\right)
=O_b\!\left(\frac1{\log Z}\right).
\tag{2.6}
\]
Thus, with
\[
\Lambda=\sum_{e\in\mathcal E}\frac1e,
\tag{2.7}
\]
one has
\[
\Lambda\longrightarrow\lambda_3.
\tag{2.8}
\]
Fix constants
\[
\frac13<\lambda_-<\lambda_3<\lambda_+<1.
\tag{2.9}
\]
For all sufficiently large $X$, $\Lambda\in(\lambda_-,\lambda_+)$. Define
\[
\theta=\frac1{b\Lambda}.
\tag{2.10}
\]
Then $\theta$ remains in a fixed compact subinterval of $(0,1)$ and
\[
\theta\Lambda=\frac1b.
\tag{2.11}
\]

We next compute the inverse-square scale. The complete-pair identity gives
\[
\sum_{p<q\in\mathcal P}\frac1{p^2q^2}
=
\frac12\left(U_X^2-\sum_{p\in\mathcal P}\frac1{p^4}\right).
\]
A dyadic PNT estimate gives
\[
U_X\sim\frac1{X\log X},
\tag{2.12}
\]
and the fourth-power term is negligible. Hence
\[
\sum_{e\in\mathcal E_{\mathrm{pair}}}\frac1{e^2}
\sim
\frac1{2X^2\log^2X}.
\tag{2.13}
\]
The target-row square mass is
\[
\sum_{e\in\mathcal E_b}\frac1{e^2}
=
\left(\sum_{r\in S_b}\frac1{r^2}\right)
\left(\sum_{q\in\mathcal B}\frac1{q^2}\right)
=O_b\!\left(\frac1{Z\log Z}\right),
\tag{2.14}
\]
which is $o(X^{-2}\log^{-2}X)$ because $Z=X^3$.

Let $(\xi_e)_{e\in\mathcal E}$ be independent Bernoulli variables with parameter $\theta$, and set
\[
S=\sum_{e\in\mathcal E}\frac{\xi_e}{e}.
\tag{2.15}
\]
Then, by \eqref{2.11},
\[
\mathbb E S=\frac1b,
\tag{2.16}
\]
and
\[
\sigma^2:=\operatorname{Var}(S)
=
\theta(1-\theta)
\sum_{e\in\mathcal E}\frac1{e^2}.
\tag{2.17}
\]
Since
\[
\theta\longrightarrow
\alpha_b:=\frac{2}{b(\log3)^2},
\tag{2.18}
\]
we obtain
\[
\sigma^2
\sim
\frac{\alpha_b(1-\alpha_b)}{2X^2\log^2X},
\qquad
\sigma^{-1}\asymp_b X\log X.
\tag{2.19}
\]

\subsection{Family size and lattice scale}

\begin{lemma}\label{lem:family-size}
As $X\to\infty$,
\[
M:=|\mathcal E|
=(1+o(1))\frac{X^6}{18\log^2X},
\tag{2.20}
\]
\[
\log L=(1+o(1))X^3=o(M),
\tag{2.21}
\]
and
\[
L\sigma\longrightarrow\infty.
\tag{2.22}
\]
\end{lemma}

\begin{proof}
The prime number theorem gives
\[
|\mathcal P|
=(1+o(1))\frac{X^3}{3\log X},
\]
so
\[
|\mathcal E_{\mathrm{pair}}|
=\binom{|\mathcal P|}{2}
=(1+o(1))\frac{X^6}{18\log^2X}.
\]
The target rows have only $O_b(Z/\log Z)=O_b(X^3/\log X)$ elements, which is negligible at this scale. This proves \eqref{2.20}.

By Chebyshev's form of the prime number theorem,
\[
\log L
=
O_b(1)+\sum_{X\le p<X^3}\log p
=(1+o(1))X^3,
\]
which gives \eqref{2.21}. Finally \eqref{2.19} and the exponential size of $L$ immediately imply \eqref{2.22}.
\end{proof}

Because every $e\in\mathcal E$ divides $L$, the random variable $LS$ is integer-valued. Thus the exact event
\[
S=\frac1b
\tag{2.23}
\]
is a single lattice atom, corresponding to $LS=L/b$. The purpose of Sections \ref{sec:fourier}--\ref{sec:gaussian} is to determine its asymptotic mass.
